\documentclass[titlepage]{article}
\usepackage{tikz}
\usepackage{xcolor}
\usepackage{graphicx}
\usepackage{grffile}
\usepackage{longtable}
\usepackage[margin=1.0in]{geometry}

\setlength{\parindent}{0cm} 
\include{r_posterior}

\newcommand{\degree}{$^{\circ}$}
\begin{document}

\begin{titlepage}
  \centering
  \vspace{4cm}
  {\huge\bfseries IceCube Fast-Response \\ Posterior Supplemental Report\par}
  \vspace{1cm}
  {\LARGE For Internal Use Only\par}
  \vfill
  {\Large Source Name: \\ \itshape\sourcename\par}
  \vspace{0.5cm}
  {\Large Observation Date(s):\\ \obsdate \par}
  \vfill
  \vspace{1cm} 
  {\Large Report Generated On:\\ \reportdate \par}
\end{titlepage}


\pagebreak
\section{Results}

{
  \centering
  {\Large Prior Zoomed}
  
  A zoom-in plot around the maximum intensity point of the skymap

  \includegraphics[width=0.9\textwidth]{\PriorSkymap}
  \clearpage
  
}
\section{Posterior Maps}

{
  \centering
  The full posterior skymap plot
  {\Large Posterior Skymap}
  \includegraphics[width=0.9\textwidth]{\PosteriorSkymap}
  \clearpage

  Zoom-in of the posterior skymap plot
  {\Large Zoomed Posterior Skymap}
  \includegraphics[width=0.9\textwidth]{\PosteriorSkymapZoom}
  \clearpage
  POsterior skymap plot zoomed in with overlaid neutrinos 
  {\Large Posterior Skymap W/ Neutrinos Overlaid}
  \includegraphics[width=0.9\textwidth]{\PosteriorSkymapZoomWNeutrinos}
  \clearpage
  
}
\pagebreak

{
  \centering
  Here we choose some pixels and plot their likelihood profile as a function of $n_s$, colored by the
  total probability of each pixel in the posterior 
  This chooses pixels from a range of probability regions by construction (since nearby points we all have very similar 
  $n_s$ profiles), but could skip over interesting pixels. 
  {\Large $N_s$ Profile Plot}
  \includegraphics[width=0.9\textwidth]{\NsProfilePlot}
  \clearpage
  
}
  



\vfill

\end{document}
